%%%%%%%%%%%%%%%%%%%%%%%%%%%%%%%%%%%%%%%%%%%%%%%%%%%%%%%%%%%%%%
%%%%%%%%%%%%%%%%%%%% MA-0840 Probabilidad %%%%%%%%%%%%%%%%%%%%
%%%%%%%%%%%%%%%%%%%%%%%%%%%%%%%%%%%%%%%%%%%%%%%%%%%%%%%%%%%%%%
\newpage
\subsection*{MA-0840 Probabilidad}
\addcontentsline{toc}{subsection}{MA-0840 Probabilidad}

\hypertarget{MA0840}{}
\begin{refsection}
\begin{table}[H]
\centering{
\begin{tabular}{|ll|}
\hline
\textbf{Año:} IV  & \textbf{Requisitos:} MA-0705 Análisis Real I \\ 
\textbf{Ciclo:} Optativa  & \textbf{Correquisitos:} Ninguno \\ 
\textbf{Tipo de curso:} Teórico & \textbf{Horas presenciales por semana:} 5 \\ 
\textbf{Créditos:} 5 & \textbf{Horas de trabajo independiente por semana:} 10 \\ \hline
\end{tabular}
}
\end{table}

\subsubsection*{Descripción del curso}

Este curso optativo introduce la teoría general de la probabilidad desde el enfoque axiomático de Kolmogorov, basado en teoría de la medida e integración. La probabilidad ha sido estudiada de manera informal durante siglos, especialmente por su incidencia en los juegos de azar; su tratamiento matemático formal se consolidó a partir del siglo XVI con Cardano y, posteriormente, con Fermat, Pascal, Huygens, Bernoulli y Laplace. En la actualidad es una de las ramas de la matemática con mayor aplicación en otras ciencias.

La probabilidad trata sobre la cuantificación de las expectativas de que un evento ocurra. El curso utiliza el sistema axiomático más extendido en la literatura moderna, apoyado en $\sigma$-álgebras, medidas de probabilidad e integración respecto a dichas medidas.

Se supone que la persona estudiante cuenta con conocimientos previos en teoría de integración sobre $\R^n$ (como los desarrollados en MA-0705 Análisis Real I), pero no se requiere experiencia previa en probabilidad. Al inicio del curso se incluye una sección de probabilidad elemental y técnicas combinatorias básicas para el cálculo de probabilidades en espacios finitos.

\subsubsection*{Objetivos}

\textbf{General:} Desarrollar los fundamentos de la teoría de la probabilidad en el marco de la teoría de la medida para modelar incertidumbre, demostrar resultados centrales y aplicarlos en problemas matemáticos e interdisciplinarios.

\textbf{Específicos:} Al finalizar el curso se espera que la persona estudiante sea capaz de:

\begin{enumerate}
\item Aplicar combinatoria elemental para resolver problemas de probabilidad en espacios finitos.
\item Identificar y utilizar los conceptos probabilísticos (eventos, $\sigma$-álgebras, variables aleatorias, esperanza, independencia) en el contexto de problemas de aplicación.
\item Demostrar resultados de probabilidad usando teoría de la medida e integración.
\item Calcular y analizar probabilidades y distribuciones mediante teoremas límite (leyes de grandes números, teorema del límite central y resultados afines).
\item Formular y analizar modelos probabilísticos sencillos mediante un proyecto que integre los contenidos del curso.
\item Comunicar argumentos probabilísticos con rigor, tanto de forma oral como escrita.
\end{enumerate}

\subsubsection*{Contenidos}

\begin{enumerate}

\item \textbf{Probabilidad elemental (3 semanas):}
\begin{enumerate}
    \item Eventos, espacios muestrales y $\sigma$-álgebras.
    \item Medidas de probabilidad; probabilidad condicional e independencia.
    \item Técnicas combinatorias para el cálculo de probabilidades.
\end{enumerate}

\item \textbf{Teoría de la medida e integración (2 semanas):}
\begin{enumerate}
    \item Conjuntos medibles, funciones medibles y clases monótonas.
    \item Integración de funciones no negativas e integrables.
    \item Producto de $\sigma$-álgebras, medidas producto, teorema de Fubini, convolución y fórmula de cambio de variables.
\end{enumerate}

\item \textbf{Teoría de probabilidad (3 semanas):}
\begin{enumerate}
    \item Espacios de probabilidad; variables aleatorias y esperanza.
    \item Funciones de distribución; $\sigma$-álgebra generada por una variable aleatoria.
    \item Momentos, varianza, regresión lineal; funciones característica y generadora.
\end{enumerate}

\item \textbf{Independencia (2 semanas):}
\begin{enumerate}
    \item Independencia de $\sigma$-álgebras; lema de Borel--Cantelli.
    \item Construcción de medidas producto; sumas de variables independientes y convolución.
    \item Proceso de Poisson.
\end{enumerate}

\item \textbf{Convergencia de variables aleatorias (3 semanas):}
\begin{enumerate}
    \item Modos de convergencia (c.s., en probabilidad, en $L^p$, en distribución).
    \item Ley débil y ley fuerte de grandes números.
    \item Convergencia en distribución; convergencia de medidas empíricas.
    \item Teorema del límite central; convergencia de Poisson.
\end{enumerate}

\item \textbf{Probabilidad condicional (2 semanas):}
\begin{enumerate}
    \item Definición de esperanza condicional (casos integrable y no negativo).
    \item Propiedades y técnicas de evaluación de la esperanza condicional.
\end{enumerate}

\end{enumerate}

\subsubsection*{Metodología}

\noindent \textbf{Lineamientos metodológicos:} La persona docente a cargo de este curso debe respetar las siguientes pautas:

\begin{enumerate}
    \item El curso combina \textbf{clases magistrales} con \textbf{sesiones de ejercicios} semanales; en estas últimas la persona estudiante presenta y discute soluciones de tareas asignadas con anticipación.
    \item Se espera participación activa: planteamiento de preguntas, discusión de ideas y comunicación rigurosa de argumentos en pizarra o equivalente.
    \item El trabajo independiente en casa es fundamental para la resolución de problemas y la asimilación de la teoría de la medida aplicada a probabilidad.
    \item Se incluye un \textbf{proyecto} desarrollado a lo largo del curso, en el que la persona estudiante investiga un tema vinculado con los contenidos, elabora un informe y realiza una presentación oral, preferiblemente en trabajo colaborativo con acompañamiento docente.
    \item Pueden asignarse lecturas complementarias (sin evaluación sumativa obligatoria) para contextualizar el desarrollo histórico y social de la probabilidad y la estadística.
    \item Cuando las autoridades sanitarias o institucionales impongan restricciones, las sesiones presenciales pueden trasladarse a modalidad virtual según indique la persona docente.
\end{enumerate}

\textbf{Sugerencias metodológicas:} Adicionalmente, se tienen las siguientes recomendaciones para la persona docente a cargo de este curso:

\begin{enumerate}
    \item Explicitar las conexiones entre los resultados de MA-0705 (medida e integración de Lebesgue) y su uso en el formalismo de Kolmogorov.
    \item Alternar demostraciones en clase con ejemplos que muestren aplicaciones en otras áreas de la ciencia.
    \item Utilizar plataformas institucionales de mediación virtual para entrega de tareas y materiales, cuando corresponda.
    \item Promover estrategias de evaluación formativa y sumativa coherentes con el marco pedagógico de la carrera: tareas periódicas, presentaciones orales, exámenes parciales y proyecto final.
    \item Fomentar el trabajo en equipo y la lectura de fuentes primarias o libros de referencia, por ejemplo \cite{AshDD99}, \cite{Bil95}, \cite{Dur19} y \cite{Gri20}.
\end{enumerate}

\nocite{AshDD99}
\nocite{Bil95}
\nocite{Dur19}
\nocite{Gri20}
\nocite{Kal21}
\nocite{LeG22}

\printbibliography[heading=subbibliography,section=\the\value{refsection}]
\end{refsection}
